\scanpage{165}
va voir que la donnée des \deleted{\unsure{$K$-}}structures \inserted{de $\underline U$ modules} sur les
$V\underset{k}{\otimes}W$ ($V,W\in\underline K^U$) suffit déjà à déterminer
un $\underline U\longrightarrow\underline U\mathbin{\widehat\otimes}\underline U$ \uline{unique} tel que ce soit défini
par lui. Pour ceci, il suffit, pour tout
voisinage $\mathcal J$, de \deleted{\unsure{se}} définir
$\underline U\longrightarrow\underline U/\mathcal J\otimes\underline U/\mathcal J$. \inserted{On peut supposer que} $\mathcal J$ est l'annulateur
d'un $V\in K^{\underline U}$, donc $\underline U/\mathcal J\inserted{={}\underline U_V}\subset L_k(V)$,
$\underline U/\mathcal J\underset{k}{\otimes}\underline U/\mathcal J\subset L_k(V)\underset{k}{\otimes}L_k(W)=L_k W\underset{k}{\otimes}V)$.

\deleted{il suffit donc} Or on \unsure{connaît|connait} bien
$\underline U\longrightarrow L_k(\unsure{V|W}\otimes V)$, cela définit ce qu'on veut
pourvu qu'on fasse la supposition

Les $\underline U$-homothéties de $V\otimes V$
$(V\in K^{\underline U})$ sont de la forme $\sum u_i\otimes v_i$,
où $u_i$ $v_i$ sont des $\underline U$-homothéties
de $V$.
\margin{\deleted{\unsure{Si $K^{\underline U}$ est|$k^{\underline U}\to$}} Si $K^{\underline U}$ est $1/2$ simple, cette condition est vérifiée d'elle-même}
% slanted

Ainsi on définit $\underline U\longrightarrow\underline U/\mathcal J\underset{k}{\otimes}\underline U/\mathcal J$.
On voit aussitôt les compatibilités
nécessaires pour que $\underline U\longrightarrow\underline U\mathbin{\widehat\otimes}\underline U$ soit
\unsure{défini|définié}, et on vérifie que cet hom.
définit la loi tensorielle cherchée.

% archive sheet number: 164
